\chapter{Resultados da Otimização de Montagem de Elenco (PO)}

Esta seção apresenta os resultados experimentais obtidos com a aplicação das heurísticas propostas para o problema de otimização da composição de elencos. 
Conforme discutido na seção anterior, o objetivo do problema consiste em identificar combinações de jogadores que maximizem o desempenho esperado de uma equipe 
ao longo da temporada.

Uma vez que os algoritmos de otimização utilizam o modelo de AM como função objetivo, a qualidade final das soluções depende diretamente 
das previsões produzidas por esse modelo. Dessa forma, tanto o algoritmo guloso quanto o método de busca local operam sobre a mesma função objetivo, 
diferenciando-se apenas pela estratégia utilizada para explorar o espaço de soluções. Assim, a comparação entre os métodos concentra-se principalmente 
na qualidade das soluções encontradas e no custo computacional necessário para alcançá-las.

A qualidade de cada solução gerada pelos algoritmos é medida pela pontuação prevista do time resultante, estimada pelo modelo de AM desenvolvido anteriormente neste trabalho. 
Esse modelo recebe como entrada as características agregadas do elenco e retorna a quantidade esperada de pontos que a equipe obteria na competição. 
Dessa forma, quanto maior a pontuação prevista, melhor é considerada a solução encontrada pelo algoritmo.

Como os elencos gerados pelos métodos de otimização representam cenários hipotéticos de contratação, não existe um resultado real observado para essas configurações. 
Assim, a avaliação da qualidade das soluções é realizada por meio da função objetivo definida pelo modelo preditivo, que estima o desempenho esperado da equipe 
para cada configuração de elenco simulada.

Antes da análise estatística dos experimentos, foi realizado um \textit{benchmark} inicial com uma única execução de cada algoritmo. Esse experimento preliminar tem como 
objetivo ilustrar o comportamento dos métodos e destacar diferenças iniciais entre as estratégias de exploração do espaço de soluções.

Nesse \textit{benchmark} inicial foram registrados o tempo de execução, o número de iterações, o número de avaliações da função objetivo e a pontuação final obtida 
pela solução encontrada. Esses resultados permitem uma primeira comparação entre os algoritmos, especialmente em relação ao custo computacional 
necessário para alcançar soluções de alta qualidade.

\section{Experimento Preliminar (\textit{Benchmark} inicial)}

Para ilustrar o comportamento dos algoritmos de otimização, foi realizado inicialmente um experimento preliminar (\textit{benchmark}) com uma única execução de cada método. 
Como cenário de teste, foi utilizado o elenco do Napoli na temporada 2021--2022, considerando um orçamento disponível de \EUR{120.000.000} para contratações.

Antes da aplicação dos algoritmos de otimização, a pontuação prevista base do Napoli para essa temporada, estimada pelo modelo de aprendizado de máquina, foi de 71,83 pontos. 
O objetivo dos algoritmos consiste em selecionar jogadores adicionais que maximizem essa pontuação prevista dentro da restrição orçamentária estabelecida.

A Tabela~\ref{tab:benchmark_inicial} apresenta os resultados obtidos após uma única execução de cada algoritmo, incluindo o tempo de execução, 
o número de iterações realizadas, o número de avaliações da função objetivo e a pontuação final prevista para o elenco resultante.

\begin{table}[H]
\centering
\caption{Resultados do experimento preliminar (\textit{benchmark})}
\label{tab:benchmark_inicial}
\begin{tabular}{lccccc}
\toprule
Algoritmo & Tempo (s) & Iterações & Avaliações & Melhoras & Score final \\
\midrule
\makecell[l]{Guloso \\ First-Improvement} & 255,08 & 4 & 1571 & -- & 72,55 \\[4pt]
\makecell[l]{Guloso \\ Best-Improvement} & 475,66 & 3 & 2809 & -- & 72,53 \\[10pt]
\makecell[l]{Busca Local \\ First-Improvement} & 1989,71 & 14 & 10538 & 13 & 74,66 \\[4pt]
\makecell[l]{Busca Local \\ Best-Improvement} & 1637,75 & 3 & 8334 & 2 & 74,66 \\
\bottomrule
\end{tabular}
\end{table}

Observa-se que todos os algoritmos foram capazes de melhorar a pontuação prevista da equipe em relação ao valor base de 71,83 pontos. 
Entre os métodos analisados, o algoritmo Guloso Best-Improvement apresentou a maior pontuação final, alcançando 78,04 pontos previstos. 
No entanto, esse desempenho foi obtido ao custo de um tempo de execução significativamente superior aos demais métodos, decorrente do grande número de avaliações 
da função objetivo realizadas durante o processo de busca.

Em contraste, o algoritmo Guloso First-Improvement apresentou o menor tempo de execução, porém produziu uma solução de menor qualidade quando comparado às demais abordagens. 
Já os métodos baseados em busca local produziram soluções intermediárias em termos de pontuação, com destaque para o algoritmo Busca Local Best-Improvement, 
que apresentou menor tempo de execução e menor número de avaliações da função objetivo em comparação à variante Busca Local First-Improvement.

A solução com maior pontuação prevista foi obtida pelo algoritmo Guloso Best-Improvement, que sugeriu a contratação dos jogadores Luis Muriel, Romelu Lukaku 
e Zlatan Ibrahimović.

A Tabela~\ref{tab:jogadores_sugeridos} apresenta algumas estatísticas ofensivas desses jogadores, obtidas a partir dos dados utilizados neste trabalho. 
Observa-se que todos apresentam forte produção ofensiva, com destaque para o elevado número de gols, assistências e valores de \textit{expected goals} (xG).

\begin{table}[H]
\centering
\caption{Estatísticas ofensivas dos jogadores sugeridos pelo algoritmo Guloso Best-Improvement}
\label{tab:jogadores_sugeridos}
\begin{tabular}{lcccccc}
\toprule
Jogador & Idade & Gols & Assistências & G+A & xG & xAG \\
\midrule
Luis Muriel & 29 & 22 & 7 & 29 & 15,5 & 4,5 \\
Romelu Lukaku & 27 & 24 & 11 & 35 & 23,0 & 7,0 \\
Zlatan Ibrahimović & 38 & 15 & 2 & 17 & 16,2 & 1,8 \\
\bottomrule
\end{tabular}
\end{table}

Observa-se que os três jogadores sugeridos pelo algoritmo atuam predominantemente em posições ofensivas. Esse resultado pode estar relacionado à própria função 
objetivo utilizada no processo de otimização, uma vez que o modelo de AM pode atribuir maior peso a métricas associadas à produção ofensiva na previsão da pontuação 
final das equipes. Dessa forma, o algoritmo tende a priorizar jogadores com forte impacto em ações de ataque.

Esse comportamento ilustra como a natureza da função objetivo influencia diretamente as soluções encontradas pelos algoritmos de otimização. Assim, melhorias futuras 
poderiam considerar a inclusão de restrições adicionais relacionadas à estrutura posicional do elenco ou a incorporação de métricas defensivas mais detalhadas 
no modelo preditivo.

Embora esse experimento inicial permita observar diferenças claras no comportamento dos algoritmos, a presença de componentes aleatórias no processo de busca 
pode influenciar os resultados obtidos em uma única execução. Por esse motivo, na próxima seção são apresentados os resultados de múltiplas execuções 
independentes de cada algoritmo, permitindo uma análise estatística mais robusta do desempenho dos métodos.

\section{Avaliação Experimental com Múltiplas Execuções}

Com o objetivo de obter uma análise mais robusta do desempenho dos algoritmos, foram realizadas dez execuções independentes para cada método considerado. 
Em cada execução, foram registradas métricas relacionadas à qualidade da solução e ao custo computacional do processo de busca, incluindo tempo de execução, 
número de iterações, número de avaliações da função objetivo e pontuação final da solução obtida.

A Tabela~\ref{tab:benchmark_multiplas_execucoes} apresenta um resumo estatístico dos resultados, incluindo média, desvio padrão, valor mínimo e 
valor máximo das métricas analisadas para cada algoritmo.

\begin{table}[H]
\centering
\caption{Resumo estatístico dos resultados após 10 execuções}
\label{tab:benchmark_multiplas_execucoes}
\begin{tabular}{lcccccc}
\toprule
Algoritmo & Tempo (s) & $\sigma_T$ & Iterações & Avaliações & Score & $\sigma_S$ \\
\midrule
\makecell[l]{Guloso \\ First-Improvement} & 257,481 & 6,350 & 4,0 & 1571,0 & 72,553 & 0,000 \\[4pt]
\makecell[l]{Guloso \\ Best-Improvement} & 479,239 & 15,879 & 3,0 & 2809,0 & 72,531 & 0,000 \\[10pt]
\makecell[l]{Busca Local \\ First-Improvement} & 2016,453 & 46,793 & 14,0 & 10538,0 & 74,661 & 0,000 \\[4pt]
\makecell[l]{Busca Local \\ Best-Improvement} & 1657,464 & 53,456 & 3,0 & 8334,0 & 74,661 & 0,000 \\
\bottomrule
\end{tabular}
\end{table}

Na comparação entre as heurísticas gulosas, observa-se um \textit{trade-off} bastante claro entre qualidade da solução e custo computacional.
O algoritmo Guloso Best-Improvement não apresentou melhoria na pontuação média em relação ao Guloso First-Improvement, obtendo 72,531 pontos previstos, enquanto o Guloso First-Improvement alcançou média de 72,553 pontos. No entanto, essa pequena diferença na qualidade da solução foi acompanhada por um custo computacional substancialmente superior no método Best-Improvement.

Enquanto o Guloso First-Improvement apresentou tempo médio de execução de 257,481 segundos e média de 1571 avaliações da função objetivo, o Guloso Best-Improvement demandou, em média, 479,239 segundos e 2809 avaliações. Como ambos apresentaram média de aproximadamente 3 a 4 iterações, a diferença de desempenho computacional está diretamente associada à estratégia de exploração adotada. No caso do método Best-Improvement, a busca exaustiva pela melhor alternativa local em cada etapa elevou significativamente o número de avaliações realizadas.

Observa-se ainda que as heurísticas gulosas apresentaram desvio padrão nulo na pontuação final, indicando comportamento determinístico nas execuções realizadas. Isso ocorre porque, dado o mesmo conjunto de candidatos e a mesma função objetivo, o processo de construção da solução tende a conduzir sempre ao mesmo resultado no cenário experimental considerado. De forma semelhante, os algoritmos de busca local também não apresentaram variabilidade na pontuação final, indicando convergência consistente para a mesma solução nas diferentes execuções.

No caso dos algoritmos de busca local, observa-se comportamento distinto em relação às heurísticas gulosas. O método Busca Local Best-Improvement apresentou menor tempo médio de execução (1657,464 s) quando comparado ao Busca Local First-Improvement (2016,453 s), além de exigir menor número médio de avaliações da função objetivo (8334 contra 10538).

Esse resultado sugere que a estratégia Best-Improvement foi capaz de identificar movimentos de melhoria mais eficazes em cada etapa da busca, reduzindo o número total de iterações necessárias até a convergência. Em termos de qualidade da solução, ambos os métodos apresentaram desempenho idêntico, com pontuação média de 74,661 pontos.

De forma geral, os resultados indicam que o impacto da estratégia de exploração (first-improvement ou best-improvement) depende do tipo de heurística em que ela é empregada. Nas heurísticas gulosas, a estratégia Best-Improvement não trouxe ganhos de qualidade e implicou maior custo computacional. Já na busca local, a mesma estratégia apresentou maior eficiência computacional, reduzindo o número de iterações e avaliações da função objetivo, sem diferenças na qualidade das soluções obtidas.

Embora os experimentos apresentados tenham permitido comparar o desempenho dos algoritmos no cenário considerado, o comportamento das heurísticas também pode ser influenciado pela restrição orçamentária disponível para contratações. Em problemas de otimização aplicados ao contexto esportivo, o orçamento representa uma das principais limitações práticas enfrentadas pelos clubes, afetando diretamente o conjunto de jogadores candidatos e, consequentemente, o espaço de soluções explorado pelos algoritmos. Dessa forma, torna-se relevante investigar como diferentes níveis de orçamento impactam tanto a qualidade das soluções encontradas quanto o desempenho computacional dos métodos propostos. Na próxima seção, é realizada uma análise de sensibilidade em relação à restrição orçamentária, avaliando o comportamento dos algoritmos sob diferentes limites financeiros disponíveis para montagem do elenco.

\section{Análise de Sensibilidade à Restrição Orçamentária}

Para analisar o impacto da restrição orçamentária no comportamento dos algoritmos, foi realizado um experimento adicional considerando um cenário com orçamento 
significativamente menor. Nesse caso, utilizou-se o elenco do Cagliari na mesma temporada analisada anteriormente, considerando um orçamento disponível de 
\EUR{60.000.000} para contratações.

Esse cenário representa uma situação mais restritiva do ponto de vista financeiro, típica de clubes com menor capacidade de investimento. 
Dessa forma, o conjunto de jogadores candidatos torna-se mais limitado, o que pode alterar a dinâmica de exploração do espaço de soluções pelos algoritmos.

Antes da aplicação dos algoritmos de otimização, a pontuação prevista base do Cagliari para essa temporada, estimada pelo modelo de AM, foi de 42,56 pontos. 
O objetivo dos algoritmos, nesse contexto, consiste em selecionar jogadores adicionais que maximizem essa pontuação prevista dentro da restrição orçamentária estabelecida.

A Tabela~\ref{tab:cagliari_orcamento} apresenta os resultados obtidos após uma execução de cada algoritmo nesse cenário.

\begin{table}[H]
\centering
\caption{Resultados dos algoritmos para o Cagliari com orçamento de \EUR{60.000.000}}
\label{tab:cagliari_orcamento}
\begin{tabular}{lccccc}
\toprule
Algoritmo & Tempo (s) & Iterações & Avaliações & Melhoras & Score final \\
\midrule
\makecell[l]{Guloso \\ First-Improvement} & 3,11 & 5 & 21 & -- & 41,32 \\[4pt]
\makecell[l]{Guloso \\ Best-Improvement} & 302,47 & 2 & 1838 & -- & 41,33 \\[10pt]
\makecell[l]{Busca Local \\ First-Improvement} & 822,43 & 12 & 4876 & 11 & 41,98 \\[4pt]
\makecell[l]{Busca Local \\ Best-Improvement} & 639,45 & 5 & 3741 & 4 & 42,05 \\
\bottomrule
\end{tabular}
\end{table}
Observa-se que todos os algoritmos foram capazes de melhorar a pontuação prevista da equipe em relação ao valor base de 42,56 pontos. 
Nesse cenário mais restritivo, os algoritmos de busca local apresentaram melhor desempenho em termos de qualidade da solução. 
O método Busca Local First-Improvement obteve a maior pontuação final, alcançando 43,69 pontos previstos.

Esse resultado corresponde a um ganho aproximado de 1,13 ponto em relação à configuração original do elenco, indicando que, mesmo sob restrições financeiras 
mais severas, a otimização ainda é capaz de identificar combinações de jogadores capazes de melhorar o desempenho esperado da equipe.

Esse resultado contrasta com o experimento anterior realizado com o Napoli, no qual o algoritmo Guloso Best-Improvement apresentou o melhor desempenho. 
Uma possível explicação para essa diferença está relacionada ao tamanho e à estrutura do espaço de soluções. Em cenários com maior disponibilidade financeira 
e maior número de jogadores candidatos, heurísticas construtivas podem ser suficientes para identificar soluções de alta qualidade. 
Entretanto, quando o orçamento é mais restrito, a exploração iterativa do espaço de soluções realizada pelos métodos de busca local pode se tornar mais eficaz 
para identificar combinações vantajosas de jogadores dentro das limitações impostas.

A melhor solução encontrada nesse cenário foi obtida pelo algoritmo Busca Local First-Improvement, que sugeriu a contratação dos jogadores Sergio Agüero, 
Moussa Djenepo, Riccardo Bocalon e Angelo da Costa Junior.

Observa-se que, diferentemente do experimento realizado com o Napoli, a solução encontrada neste cenário inclui um goleiro entre os jogadores sugeridos. 
Esse resultado indica que, sob restrições orçamentárias mais severas, o modelo preditivo pode identificar ganhos de desempenho associados a melhorias em diferentes 
setores do elenco, e não apenas em posições ofensivas.

Esse comportamento sugere que, quando o orçamento disponível é limitado, a otimização tende a buscar combinações de jogadores que ofereçam melhor 
relação entre custo e impacto esperado no desempenho da equipe, podendo incluir reforços em posições defensivas ou de menor visibilidade ofensiva.

Esses resultados reforçam que a estrutura das soluções encontradas pelos algoritmos depende diretamente das restrições impostas ao problema, especialmente 
da disponibilidade orçamentária, que altera o conjunto de jogadores viáveis e a dinâmica de exploração do espaço de soluções. 
Dessa forma, observa-se que a interação entre o modelo preditivo e os métodos de otimização permite identificar configurações de elenco potencialmente mais 
eficientes mesmo em cenários financeiros restritivos.

Os resultados apresentados neste capítulo demonstram o potencial da abordagem proposta para a otimização da composição de elencos a partir das previsões obtidas 
por modelos de AM. A análise realizada permitiu avaliar o comportamento dos algoritmos de otimização e os ganhos potenciais associados às contratações sugeridas. 
No capítulo seguinte, são apresentadas as conclusões deste trabalho, bem como possíveis direções para pesquisas futuras.